\section{Discussion}
\label{sec:limitations}

The minimal worlds separate two demands on an abstract world model. Its representation must retain
the variables on which a decision depends. Its prediction must preserve the expectations that the
planner will evaluate. A non-collapsed representation can fail the first demand by concentrating on
predictable nuisance. A correct conditional-mean model can fail the second when a planner applies
nonlinear costs or composes nonlinear dynamics at that mean.

The redcube comparisons make these demands visible at different stages: requested-goal movement,
position accessible in real observations, and behavior learned from imagined futures. The matched
nuisance study already shows a reproducible change in accessible state under combined fresh noise
and wind. Its high prediction loss separates the measured finite-budget failure from the idealized
nuisance optimum. Likewise, a practical JEPA/flow control gap would compare complete learned
packages; the exact variance calculation is the narrower test that isolates point compression.

The empirical scope is small visual domains, specified models and training budgets, and the
reported checkpoint and policy replications. The elementary examples establish limitations of
mean-state planning, not a requirement that every useful world model generate samples. Expected-cost
or value-equivalent models can preserve the needed decision quantities directly. General-purpose
world models need a principle for choosing both the state to preserve and the predictive information
that their intended decisions require.
